% =====================================================================
%  CHƯƠNG 9. KẾT LUẬN VÀ HƯỚNG PHÁT TRIỂN         | Phụ trách: Cả 2
%  Ngân sách: mục tiêu 3 trang, trần 1500 từ (.tex)
%  Số nguồn gốc:
%    mAP@0,5 detector:    src/ml/experiments/a1_dedup_eval.csv (0,983)
%    OCR topkek 56,6%:    src/ml/experiments.csv (2026-08-29)
%    OCR vn_plate 80,2%:  docs/report/tuan-07.md mục 1
%    Loại xe 98,66%:      src/ml/experiments/type_resnet18_history.json (ResNet18)
%    Kiểu dáng 90,07%:    src/ml/experiments.csv (MobileNetV3-Small, model triển khai)
%    Pi5 0,464s/lượt:     src/ml/experiments/pi5_anh_that/tong_hop.csv (2026-09-13)
%  INT8: CHƯA thực hiện — chỉ nêu ở hướng phát triển.
% =====================================================================
\chapter{Kết Luận và Hướng Phát Triển}
\label{ch:ketluan}

\section{Tóm Tắt Đóng Góp}

Đề tài xây dựng một hệ thống quản lý bãi giữ xe thông minh với bốn đóng góp chính: (1)
pipeline nhận diện năm thành phần (Mục~\ref{sec:donggop}) dùng chung giữa máy tính và thiết
bị biên qua mô hình ONNX, không biên dịch lại; (2) mở bảng ký tự
OCR lên 37 ký tự và sửa luồng tiền xử lý lẫn suy luận để đọc đúng ký tự ``Đ'' của biển xe
máy điện; (3) backend nghiệp vụ đầy đủ với phiên gửi xe, đối chiếu chống gian lận theo hash
và Levenshtein, tính phí theo nhóm xe, phân quyền, tự xóa dữ liệu theo Luật Bảo vệ dữ liệu
cá nhân, cùng dashboard tra cứu, thống kê và xuất CSV; (4) kiến trúc hai đường thu nhận
(worker biên và kiosk máy tính) cùng đổ về một backend, đạt dưới hai giây mỗi lượt xe khi đo
trực tiếp trên Raspberry Pi~5.

\section{Trả Lời Câu Hỏi Nghiên Cứu}

Bốn câu hỏi nghiên cứu ở Mục~\ref{sec:cauhoi} được trả lời như sau, trong phạm vi thử nghiệm
của đề tài.

\textbf{CH1 (độ tin cậy).} OCR đọc đúng cả biển 91,9\% trên tập ảnh camera cổng a1 và loại xe
đạt 98,66\%. Với CER 0,015 trên a1, phần lớn biển sai chỉ lệch một ký tự; thuật toán đối
chiếu tự sửa khi chuỗi lệch không quá một ký tự so với phiên trong bãi, phần còn lại đưa danh
sách gợi ý cho nhân viên, nên sai sót của mô hình không chặn nghiệp vụ.

\textbf{CH2 (tối ưu tài nguyên).} Cấu hình thực thi quan trọng ngang việc chọn kiến trúc: đặt
4 luồng và dùng định dạng gốc cho bước định vị xe giảm thời gian một lượt còn 495 ms, thay
kiểu dáng sang MobileNetV3-Small hạ phân vị 90 và chạy mát hơn. Độ trễ đầu cuối trung vị
0,464 s mỗi lượt, dưới hai giây, với khoảng 3,1 J mỗi lượt.

\textbf{CH3 (phương pháp dữ liệu).} Các bộ dữ liệu công khai đủ để huấn luyện, nhưng kiểm
soát chất lượng là điều kiện tiên quyết: đề tài phát hiện và khắc phục ba lỗi học lệch do dữ
liệu (nhãn biển hai dòng chia sai, ký tự ``Đ'' bị xóa khỏi nhãn, mô hình loại xe học tắt theo
nguồn ảnh), và dùng tập kiểm tra gán nhãn thủ công độc lập để phản ánh khách quan năng lực mô
hình.

\textbf{CH4 (bảo mật và riêng tư).} Chuỗi biển và ảnh được mã hóa đối xứng, tra cứu qua giá
trị băm có khóa mà không cần giải mã; phân quyền theo vai, nhật ký thao tác nhạy cảm và tự
động xóa dữ liệu phiên sau 30 ngày đáp ứng các yêu cầu cốt lõi của Luật Bảo vệ dữ liệu cá
nhân.

\section{Đối Chiếu Mục Tiêu Đề Cương}

Bảng~\ref{tab:muctieu} đối chiếu năm mục tiêu cụ thể và hai tiêu chí thành công trong đề cương
với trạng thái thực tế của đề tài.

\begin{table}[htbp]
\centering
\caption{Đối chiếu mục tiêu đề cương với kết quả đạt được.}

\label{tab:muctieu}
\begin{tabular}{p{5cm}p{2cm}p{6cm}}
\toprule
\textbf{Mục tiêu / Tiêu chí} & \textbf{Trạng thái} & \textbf{Ghi chú} \\
\midrule
MT1: Nhận diện biển số (biển 1 hàng và 2 hàng) & Đạt & YOLOv8n mAP@0,5 = 0,989 (0,983 sau khử trùng lặp; Bảng~\ref{tab:det-dedup}, Ch~\ref{ch:phathien}); CRNN đọc đúng cả biển 91,9\% trên 172 biển camera cổng a1, biển 2 dòng không kém biển 1 dòng (Ch~\ref{ch:ocr}). \\
MT2: Phân loại loại xe và nhận màu nền biển & Đạt một phần & Loại xe ResNet18 accuracy 98,66\%; kiểu dáng MobileNetV3-Small 90,07\% (Ch~\ref{ch:phanloai}); màu biển HSV đúng 178/178 trên nhãn trắng/đỏ, vàng/xanh chưa có nhãn. \\
MT3: CSDL, logic vào ra, đối chiếu, tính phí & Đạt & Backend đầy đủ: phiên, đối chiếu hash và Levenshtein, phí flat/block/grace/daily cap (Ch~\ref{ch:hethong}). \\
MT4: Dashboard quản lý và thống kê & Đạt & Màn cổng, tra cứu phiên, thống kê, xuất CSV, phân quyền (Ch~\ref{ch:hethong}). \\
MT5: Tối ưu và triển khai trên thiết bị biên & Đạt & Pipeline xuất ONNX, chạy 4 luồng; đo trực tiếp trên Raspberry Pi 5, trung vị 0,464 s mỗi lượt xe, khớp máy dev 50/50 ảnh (Ch~\ref{ch:trienkhai}); chưa thực hiện lượng tử hóa INT8. \\
\midrule
TC1: Độ chính xác đọc biển $>$\SI{90}{\percent} điều kiện tiêu chuẩn & Đạt trên tập a1 & 91,9\% đọc đúng cả biển trên 172 biển camera cổng a1 (KTC 95\% từ 86,8\% đến 95,1\%); 80,2\% trên vn\_plate và 56,6\% trên topkek (ảnh độ phân giải thấp) (Ch~\ref{ch:ocr}, Ch~\ref{ch:trienkhai}). \\
TC2: Độ trễ $<$ \SI{2}{giây}/xe trên Raspberry Pi 5 & Đạt & Đo trực tiếp trên Raspberry Pi 5 thật: trung vị 0,464 s mỗi lượt xe qua API, p90 dưới 0,5 s (Mục~\ref{sec:trienkhai-e2e}, Ch~\ref{ch:trienkhai}). \\
\bottomrule
\end{tabular}
\end{table}

Về sản phẩm bàn giao, phần mềm end-to-end cùng dashboard đã hoàn thiện và mã nguồn được quản
lý trên hệ thống kiểm soát phiên bản; dữ liệu gồm các bộ công khai kết hợp tập vn\_plate 96
biển gán nhãn thủ công. Video demo chưa được sản xuất trong khuôn khổ đề tài.

\section{Ứng Dụng Thực Tiễn}

Hệ thống phù hợp với bãi xe trường học, ký túc xá (tính phí theo nhóm xe, phân quyền nhân
viên cổng), chung cư (truy vấn lịch sử phiên, xuất CSV đối chiếu cư dân) và thương mại (thống
kê vận hành theo giờ và ngày). Kiến trúc mở rộng theo chiều ngang bằng cách thêm worker biên
đăng ký vào cùng một backend, và theo chiều dọc bằng cách thay detector hoặc OCR chỉ qua cập
nhật đường dẫn file ONNX, nhờ pipeline đã tách phần ONNX Runtime khỏi logic nghiệp vụ.

\section{Giới Hạn}

\textbf{Phạm vi dữ liệu.} Dữ liệu huấn luyện chủ yếu ban ngày ánh sáng đủ, chưa có số liệu
định lượng cho ban đêm, ngược sáng hoặc biển cũ mờ; tập A1 sau khử trùng lặp còn 342 ảnh;
chưa có tập test ảnh toàn cảnh cổng có nhãn chuỗi ký tự để đo accuracy toàn luồng; lớp xe
buýt chưa có mô hình riêng. Đề tài chưa tự thu thập ảnh bãi xe; phần tự thực hiện giới hạn ở
gán nhãn thủ công hai tập kiểm tra OCR độc lập.

\textbf{Domain gap kiểu dáng.} Model kiểu dáng huấn luyện chủ yếu trên ảnh đại lý của B5, nên
accuracy 90,07\% có thể không phản ánh đúng hiệu năng trên ảnh camera cổng thật với góc chụp
đa dạng và nền phức tạp.

\textbf{Màu nền biển mới đánh giá một phần.} Mô đun màu HSV cộng CLAHE phân loại đúng 178/178
crop trắng và đỏ, nhưng vàng và xanh chưa có nhãn nên chưa đủ bốn màu, trong khi màu biển là
tín hiệu phân nhóm loại phí.

\textbf{Hiệu năng thiết bị biên.} Phép đo trên Pi~5 thực hiện trên một thiết bị trong phòng
thí nghiệm có tản nhiệt chủ động; độ ổn định dài hạn và throttle nhiệt khi vận hành liên tục
ngoài trời chưa được kiểm chứng. Lượng tử hóa INT8 chưa thực hiện.

\section{Hướng Phát Triển}
\label{sec:huongphattrien}

\textbf{Ngắn hạn.} Thu thập ảnh điều kiện khó (ban đêm, ngược sáng, biển cũ) và xây tập test
ảnh toàn cảnh có nhãn chuỗi ký tự để đo accuracy end-to-end lần đầu. Thực hiện lượng tử hóa
INT8 tĩnh có tập hiệu chuẩn trên ONNX Runtime biên dịch với nhân tối ưu cho ARM, vì lượng tử
hóa động không rút ngắn độ trễ khi thiếu nhân INT8 vector hóa; đồng thời đo độ ổn định khi
worker chạy liên tục nhiều giờ. Bổ sung nhãn vàng và xanh cho tập test màu để đo đủ bốn màu và điều
chỉnh ngưỡng HSV theo ánh sáng thực tế.

\textbf{Trung hạn.} Nâng detector lên YOLO26n để tận dụng kiến trúc không cần NMS và bỏ DFL,
giảm độ trễ trên CPU. Bổ sung hàng đợi offline cho worker biên (lưu tạm và gửi lại khi phục
hồi kết nối). Tích hợp thanh toán điện tử (VNPay, MoMo) vào service tính phí đã tách độc lập.

\textbf{Dài hạn.} Mở rộng từ mô hình cổng dừng chờ sang nhận dạng dòng xe di chuyển tốc độ
cao (free-flow) với camera tốc độ cao, detector mạnh hơn và tracking qua nhiều khung; triển
khai đa cổng đa camera với một backend trung tâm kèm phân tích lưu lượng thời gian thực để
quản lý bãi xe lớn theo cụm.
